\documentclass[11pt]{article}

\usepackage[T1]{fontenc}
\usepackage{lmodern}

\usepackage[margin=1in]{geometry}
\usepackage{amsmath,amssymb,amsthm}
\usepackage{booktabs}
\usepackage{microtype}
\usepackage{xcolor}
\pagecolor{white}
\usepackage{eso-pic}
\AddToShipoutPictureBG{\AtPageLowerLeft{\color{white}\rule{\paperwidth}{\paperheight}}}
\definecolor{accent}{HTML}{142a5e}

\newtheorem{theorem}{Theorem}

\title{\textcolor{accent}{Constraint-First Physics}\\ a falsifiable route to the theory of everything}
\author{Rob Mitchell}
\date{}

\begin{document}
\maketitle

\begin{abstract}
String theory's central obstacle is not that it predicts too little --- it is that it
appears to predict \emph{nothing} testable, because the landscape of $\sim10^{500}$
vacua cannot be searched or selected. This note states the result that reframes that
obstacle as the answer: \textbf{the theory of everything predicts nothing exactly,
except what it forbids (constraints) and what it is handed (labels)}. For the flux
sector this is now a \emph{proven theorem} --- the only exact structure is the D3
tadpole, and no exact, symmetry-respecting formula can single out the special vacuum. The
frontier is the Standard Model: its exact structure is also constraint + label, and
the open question --- \emph{are the 19 free parameters unpredicable in any
ultraviolet completion?} --- is a concrete, falsifiable target that a partner can pick
up and settle either way. Everything quoted here is machine-verified.
\end{abstract}

\section{The obstacle, restated}

Two programs have tried to make the string landscape navigable. The \emph{search}
program asks whether we can learn the map from flux to observables and scan our way to
our universe. The \emph{constraint} program (swampland, bootstrap) asks which
consistency conditions eliminate vacua. The difference between them is one measurable
question: does any observable carry a \emph{fingerprint} --- strong information about
the flux --- that a search could exploit?

Measured on $174$ real supersymmetric flux vacua, the answer is \emph{no}. Every
low-energy observable (string coupling, gauge couplings, mass-spectrum shape, axions,
dilaton, cosmological constant) is statistically uniform in flux; exactly one quantity
is exactly structured --- the D3 tadpole, the constraint itself. The fingerprint the
search program needs is absent.

\section{What we proved: the theorem tower}

The constraint-first claim began as that measurement and is now a tower of six
theorems, each verified symbolically at zero tolerance, plus a machine-checked Lean arm
(\texttt{pacland-proof/}) that now machine-checks the FFT slices at degrees 1--4 (no
linear, only the tadpole, no cubic, only $N^2$). One command reproduces the whole tower
(\texttt{scripts/validate\_all\_theorems.py}).

\begin{center}
\begin{tabular}{lll}
\toprule
theorem & statement & grade \\
\midrule
Mitchell's Uniqueness & the tadpole $N=F^\top JH$ is the \emph{unique} bilinear symplectic invariant & 6/6 \\
Constraint-First & the tadpole generates the \emph{entire} invariant ring (Weyl 1939, k=2) & 5/5 \\
Mitchell's Completeness & $N$ is the \emph{unique} invariant; finitely many orbits per tadpole over $\mathbb{Z}$ & 6+2 \\
Full-duality & holds for $\mathrm{Sp}(b_3)\times\mathrm{SL}(2,\mathbb{Z})$, not just $\mathrm{Sp}$ & 5/5 \\
Complete Exact Structure & the ToE's exact content = constraints + labels, nothing else & assembly \\
\bottomrule
\end{tabular}
\end{center}

The headline consequence: \textbf{the only exact, symmetry-respecting selector in the
flux landscape is the tadpole, and it is bounded.} No exact polynomial,
duality-invariant formula can \emph{select} a vacuum; the constraint $N\le L$ can only
\emph{eliminate}. This is a no-go on exact algebraic selectors --- it does \emph{not}
rule out statistical or approximate methods, which the search program can still use.
The claim is falsifiable and sharp: \emph{a single
moduli-independent, duality-invariant flux quantity that is not a function of $F^\top
JH$ breaks the whole tower.} Nobody has to take it on faith.

\section{The reframing: the ToE's exact content}

\begin{theorem}[Complete Exact Structure]
Every exact (moduli-independent, duality-invariant) observable in a flux
compactification is either a \textbf{constraint} (the tadpole, quantization, anomalies)
or a \textbf{label} (the topology $h^{11},h^{21},c_2,\kappa$). Nothing else.
\end{theorem}

In plain terms: the theory predicts nothing exactly except what it \emph{forbids} and
what it is \emph{handed}. The quantities a theory of everything is conventionally
expected to predict --- the couplings, the mass ratios, the cosmological constant ---
are exactly the things that are \emph{not} exact, and hence not predictable. That is
not a bug in the framework; it is the answer.

\section{The frontier: the Standard Model}

The theorem's open question was \emph{are all exact structures constraints?} The
Standard Model is where that bites, and the answer so far is: \emph{yes, with the same
two-type decomposition}.

\begin{center}
\begin{tabular}{lll}
\toprule
type & flux (proven) & Standard Model \\
\midrule
constraint & tadpole $N\le L$ & anomaly cancellation, charge quantization \\
label & $h^{11},h^{21},c_2,\kappa$ & gauge group $\mathrm{SU}(3)\!\times\!\mathrm{SU}(2)\!\times\!\mathrm{U}(1)$; $N_{\rm gen}=3$ \\
parameter & $W_0$, $g_s$, couplings, $\Lambda$ & the \textbf{19 free couplings} \\
\bottomrule
\end{tabular}
\end{center}

Two exact results are already established (both machine-verified):

\begin{enumerate}
\item \textbf{Anomaly cancellation is a constraint, not a prediction.} All five SM
gauge/gravitational anomaly coefficients vanish exactly
(\texttt{validate\_sm\_anomaly.py}, 5/5). The hypercharges
$(1/6,\,2/3,\,-1/3,\,-1/2,\,-1)$ are \emph{forced} by consistency, not chosen.
\item \textbf{The hypercharge is unique.} Given the SM matter content, the anomaly
conditions determine the hypercharges uniquely up to overall normalization and the
$u\leftrightarrow d$ swap (\texttt{validate\_sm\_hypercharge\_uniqueness.py}, 4/4).
This is the gauge-sector analog of Mitchell's Uniqueness --- the \emph{constraint} is
unique.
\end{enumerate}

\section{The open target --- and why it is worth picking up}

The $19$ free parameters ($3$ gauge, $9$ masses, $4$ CKM, $2$ Higgs, $1$ $\theta_{\rm
QCD}$) are, within the Standard Model, free inputs \emph{by definition} --- that is a
trivial counting fact, not a result. The non-trivial claim is:

\begin{center}
\fbox{\parbox{0.9\textwidth}{\textbf{Conjecture (refined).} A theory of everything
predicts \emph{relations}, never \emph{values}. The known ultraviolet completions do
relate the $19$ --- gauge-coupling unification, Yukawa equalities --- but none derives
the \emph{values} correctly. Measured: no completion derives more than \textbf{2} of the
$19$ correctly (only gauge unification survives; minimal $\mathrm{SU}(5)$'s five
relations and $\mathrm{SO}(10)$'s twelve are numerically wrong).}}
\end{center}

This is the genuinely open, genuinely falsifiable frontier. The flux theorem was
provable only because flux has a clean duality group ($\mathrm{Sp}$) with a known
invariant theory (Weyl's First Fundamental Theorem). The Standard Model's couplings have
no such clean group action, so the conjecture is exactly that --- and it may be
\emph{false}: a GUT or string construction that derives a \emph{value} correctly (beyond
the two surviving relations) would be a counterexample.

\textbf{Why a partner should care.} This is a program with a success criterion, not a
slogan. Either direction is a publishable, decisive result:
\begin{itemize}
\item \emph{Confirmation:} exhibit the invariant-theory structure (the ``duality group +
FFT'' analog for the gauge sector) that proves no completion can fix the values of the
$19$ parameters --- completing the theorem in a second sector.
\item \emph{Falsification:} exhibit one ultraviolet relation among the $19$ that is
\emph{numerically correct} beyond the two surviving ones --- which would itself be the
kind of concrete prediction string theory has been accused of lacking.
\end{itemize}
Either outcome moves unification forward. The tooling is in place and open
(\texttt{github.com/robmitchellengineering-hub/pacland}); the six arXiv papers are
ready; the verification discipline (a probably-approximately-correct certificate, the
Mitchell Verification Deficit, the Rule of Three) is the same instrument used to keep
the flux result honest, and it applies unchanged.

\section{What is proven vs. what is open}

\begin{center}
\begin{tabular}{ll}
\toprule
claim & status \\
\midrule
flux: the only exact structure is the constraint & \textbf{proven} (6 theorems + Lean degrees 1--4) \\
SM: the constraint (anomaly $\to$ hypercharge) is unique & \textbf{proven} (4/4 exact) \\
SM: the $19$ parameters are free inputs & \textbf{counted} (trivial) \\
SM: the $19$ parameters are unpredicable in any ToE & \textbf{open, falsifiable} \\
\bottomrule
\end{tabular}
\end{center}

The honest status of the program: the gravity sector's exact structure is a theorem;
the gauge sector's exact structure is a theorem; the one sentence left --- \emph{``in
every sector, the theory predicts nothing except constraints and labels''} --- is proven
in one sector and a falsifiable conjecture in the other. Closing it is the next
theorem, and it has a clear yes-or-no test.

\end{document}
